% =====================================================================
%  Chương 4 — Hạng mục 4: hiện thực Keras và đối chuẩn PyTorch với Keras
%  Mã nguồn: a06/models_keras.py, a06/bench.py · Notebook: 05_keras.ipynb, 06_doi_chuan.ipynb
% =====================================================================
\chapter{Keras và đối chuẩn hai framework}
\label{ch:keras}

Chương này viết lại bốn mô hình bằng Keras~\cite{chollet2015} (backend TensorFlow~\cite{abadi2016}), huấn luyện trên
cùng dữ liệu, rồi so 16 mô hình theo chất lượng, kích thước và tốc độ. Mọi chỉ số của cả hai framework được tính lại bằng
cùng một hàm trong \texttt{a06/bench.py}, từ dự báo đã lưu, để phép so không lệch vì cách tính.

\section{Bốn mô hình bằng Keras}

Kiến trúc giữ đúng như bản PyTorch: đầu vào \texttt{Input(T, 8)}, một tầng hồi quy 64 đơn vị (một trong
\texttt{SimpleRNN}, \texttt{LSTM}, \texttt{GRU} và \texttt{Bidirectional(LSTM)}), rồi hai tầng \texttt{Dense(32, relu)}
và \texttt{Dense(1)}. Tối ưu cũng giống:
AdamW với tốc độ học $10^{-3}$, weight decay $10^{-4}$, \texttt{clipnorm=1.0}.

\maNguon{models_keras__build_keras}{Bốn mô hình Keras dựng từ một hàm}

Có một khác biệt nhỏ về clipping. \texttt{clipnorm} của Keras cắt chuẩn của \emph{từng} tensor gradient, còn
\texttt{clip\_grad\_norm\_} của PyTorch cắt theo chuẩn gộp của mọi gradient. Với ngưỡng 1 và mô hình nhỏ như ở đây, cả hai
hiếm khi được kích hoạt sau vài epoch đầu, nên bài giữ mặc định của mỗi framework.

Keras không có tham số \texttt{pos\_weight}. Cách tương đương là trọng số mẫu: mẫu dương nhận trọng số bằng số âm chia số
dương, mẫu âm nhận 1. Trọng số được truyền cho cả tập train lẫn tập validation, để val loss mà \texttt{EarlyStopping}
theo dõi đo đúng đại lượng như val loss của PyTorch.

\maNguon{models_keras__sample_weights}{Trọng số mẫu thay cho \texttt{pos\_weight}}

\maNguon{models_keras__fit_keras}{Huấn luyện một mô hình Keras với cùng lô và cùng dừng sớm}

\subsection{Chênh lệch số tham số}

Cùng kiến trúc nhưng số tham số không bằng nhau (Hình~\ref{fig:size}a). PyTorch lưu hai vector bias cho mỗi khối cổng
($b_{ih}$ và $b_{hh}$), Keras chỉ lưu một. Vì vậy Keras ít hơn đúng $H = 64$ tham số với Simple RNN
(\SL{bench/cmp/amzn_rnn/param_diff}), $4H$ với LSTM (\SL{bench/cmp/amzn_lstm/param_diff}) và $8H$ với BiLSTM
(\SL{bench/cmp/amzn_bilstm/param_diff}). GRU là ngoại lệ: Keras mặc định \texttt{reset\_after=True}, áp cổng reset sau
phép nhân ma trận giống PyTorch, nên cần hai vector bias và số tham số bằng nhau
(\SL{bench/models/amzn_torch_gru/params}). Bài giữ nguyên chênh lệch này thay vì ép bằng, vì nó không đổi khả năng biểu
diễn: tổng hai bias chỉ là một bias.

\begin{figure}[H]
  \centering
  \hinh{f4_size}
  \caption{Số tham số và kích thước tệp mô hình của hai framework}
  \label{fig:size}
\end{figure}

Kích thước tệp chênh nhiều hơn số tham số. Tệp \texttt{.pth} chỉ chứa \texttt{state\_dict}, còn tệp \texttt{.keras}
mặc định lưu thêm cấu hình mô hình và trạng thái của optimizer (hai moment của AdamW cho mỗi trọng số), nên nặng gấp
khoảng ba lần: \SL{bench/models/amzn_keras_lstm/size_kb} KB so với \SL{bench/models/amzn_torch_lstm/size_kb} KB với LSTM.
Nhờ vậy tệp Keras huấn luyện tiếp được ngay, nhưng khi chỉ cần phục vụ thì phần optimizer là thừa.

\section{Huấn luyện song song trên hai thiết bị}

TensorFlow từ bản 2.11 không còn hỗ trợ GPU trên Windows, nên Keras chạy trên CPU (16 luồng) trong một tiến trình riêng,
\emph{cùng lúc} với notebook PyTorch dùng GPU. Hai bên không phải chờ nhau, và CPU 32 luồng của máy được chia đôi cho hai
tiến trình. Hình~\ref{fig:loss-keras-amzn} và~\ref{fig:loss-keras-kkbox} là đường loss của bốn mô hình Keras; hình dạng
giống bản PyTorch ở Chương~\ref{ch:pytorch}: trên AMZN val loss đứng yên ngay từ đầu, trên KKBox cả hai đường cùng giảm.

\begin{figure}[H]
  \centering
  \hinh{f4_loss_keras_amzn}
  \caption{Đường loss của bốn mô hình Keras trên AMZN (liền: train, đứt: val)}
  \label{fig:loss-keras-amzn}
\end{figure}

\begin{figure}[H]
  \centering
  \hinh{f4_loss_keras_kkbox}
  \caption{Đường loss của bốn mô hình Keras trên KKBox (liền: train, đứt: val)}
  \label{fig:loss-keras-kkbox}
\end{figure}

Thời gian huấn luyện phản ánh chủ yếu thiết bị, không phản ánh framework. Trên KKBox, một epoch BiLSTM mất
\SL{bench/models/kkbox_keras_bilstm/epoch_s} giây với Keras trên CPU và \SL{bench/models/kkbox_torch_bilstm/epoch_s} giây
với PyTorch trên GPU. Vì vậy phần so tốc độ bên dưới đo suy luận trên cùng một thiết bị.

\section{Đo độ trễ suy luận}

Để so công bằng, độ trễ của cả 16 mô hình được đo trên CPU với đúng một luồng. Mỗi mô hình chạy khởi động 20 lần, rồi đo
200 lần với lô một mẫu và lấy trung vị (độ trễ), và đo 30 lần với lô 256 mẫu (thông lượng). Mô hình Keras chạy qua
\texttt{tf.function}, tức đồ thị đã biên dịch, là cách dùng khi phục vụ thật; gọi \texttt{model(x)} ở chế độ eager chậm
hơn khoảng một trăm lần vì vòng lặp theo thời gian chạy từng phép tính một.

\maNguon{bench___median_time}{Đo thời gian: khởi động rồi lấy trung vị}
\maNguon{bench__latency}{Độ trễ với lô 1 và thông lượng với lô 256}

Khi viết notebook này, lần đo đầu tiên cho PyTorch chậm bất thường, 1,5 đến 5 ms mỗi mẫu. Nguyên nhân là mô hình PyTorch và
Keras được đo xen kẽ nhau: sau mỗi lần gọi, các luồng của TensorFlow tiếp tục quay chờ việc một lúc và giành mất lõi CPU
duy nhất của PyTorch. Đo hết PyTorch rồi mới đo Keras thì kết quả ổn định. Đây là một ví dụ nhỏ cho thấy đối chuẩn hiệu
năng dễ sai như thế nào nếu không kiểm soát môi trường.

\section{Bảng đối chuẩn 16 mô hình}

\begin{table}[H]
  \centering\scriptsize
  \setlength{\tabcolsep}{4pt}
  \caption{Đối chuẩn trên AMZN: 8 mô hình và mốc ngây thơ}
  \label{tab:bench-amzn}
  \input{generated/tab_bench_amzn}
\end{table}

\begin{table}[H]
  \centering\scriptsize
  \setlength{\tabcolsep}{4pt}
  \caption{Đối chuẩn trên KKBox: 8 mô hình}
  \label{tab:bench-kkbox}
  \input{generated/tab_bench_kkbox}
\end{table}

\begin{figure}[H]
  \centering
  \hinh{f4_scatter}
  \caption{Chất lượng theo độ trễ suy luận của 16 mô hình}
  \label{fig:scatter}
\end{figure}

Bảng~\ref{tab:bench-amzn}, Bảng~\ref{tab:bench-kkbox} và Hình~\ref{fig:scatter} cho bốn kết luận.

\textbf{Chất lượng gần như không phụ thuộc framework.} Trên KKBox, F1 của cùng một kiến trúc ở hai framework chênh nhau
từ \SL{bench/cmp/kkbox_rnn/diffPct} (Simple RNN) đến \SL{bench/cmp/kkbox_gru/diffPct} (GRU) điểm phần trăm. Trên AMZN, RMSE chênh không
quá \SL{bench/cmp/amzn_bilstm/diff} USD, trừ Simple RNN của Keras kém hơn \SL{bench/cmp/amzn_rnn/diff} USD vì dừng ở
epoch \SL{bench/models/amzn_keras_rnn/best_epoch}. Mỗi mô hình chỉ chạy một seed, nên những chênh lệch cỡ này nằm trong mức
dao động của khởi tạo ngẫu nhiên, không đủ để nói framework nào học tốt hơn.

\textbf{PyTorch nhanh hơn với từng mẫu một.} Độ trễ trung vị của 8 mô hình PyTorch là \SL{bench/lat_median/torch} ms,
của Keras là \SL{bench/lat_median/keras} ms. Với lô một mẫu, chi phí gọi hàm của framework chiếm phần lớn thời gian, và
lời gọi eager của PyTorch nhẹ hơn lời gọi \texttt{tf.function}. BiLSTM chênh nhiều nhất: Keras chậm hơn
\SL{bench/cmp/amzn_bilstm/lat_ratio} lần.

\textbf{Với lô lớn, khoảng cách hẹp lại và có chỗ đảo chiều.} Ở lô 256, Keras Simple RNN xử lý
\SL{bench/models/kkbox_keras_rnn/throughput} mẫu/giây trên KKBox, nhiều hơn PyTorch
(\SL{bench/models/kkbox_torch_rnn/throughput}); Keras GRU trên KKBox cũng nhanh hơn
(\SL{bench/models/kkbox_keras_gru/throughput} so với \SL{bench/models/kkbox_torch_gru/throughput}). Với LSTM và BiLSTM,
PyTorch vẫn nhanh hơn, từ \SL{bench/cmp/kkbox_bilstm/thr_ratio} đến \SL{bench/cmp/kkbox_lstm/thr_ratio} lần. Chọn framework theo tốc độ cần biết trước ứng dụng phục vụ từng
yêu cầu một hay theo lô.

\textbf{Mô hình ``tốt nhất'' khác nhau tuỳ tiêu chí.} Theo val loss (tiêu chí dùng để chọn mô hình cho Web App), mô hình
tốt nhất trên AMZN là \SL{bench/best/amzn_torchLabel} (PyTorch) và \SL{bench/best/amzn_kerasLabel} (Keras), trên
KKBox là \SL{bench/best/kkbox_torchLabel} (PyTorch) và \SL{bench/best/kkbox_kerasLabel} (Keras). Chỉ số test tốt nhất
lại thuộc về \SL{bench/best_test/amznLabel} trên AMZN và \SL{bench/best_test/kkboxLabel} trên KKBox. Khi các mô hình gần bằng nhau như vậy, thứ hạng trên test gần
như là ngẫu nhiên, và mô hình rẻ nhất đủ tốt là lựa chọn hợp lý.

\section{Giới hạn của mô hình dự báo giá}

Trên AMZN, không mô hình nào trong 16 mô hình thắng mốc ngây thơ (dòng cuối Bảng~\ref{tab:bench-amzn}). Đây không phải
lỗi cài đặt. Nó đến từ ba điều đã thấy trong dữ liệu:

\begin{itemize}
  \item \textbf{Return gần như không có tự tương quan} (Chương~\ref{ch:dulieu}): quá khứ 30 phiên chứa rất ít thông tin
        về return ngày mai. Dự báo tốt nhất theo sai số bình phương khi đó là một hằng số gần 0, đúng là điều mô hình
        học được (val loss đứng yên từ epoch đầu).
  \item \textbf{Trễ pha tại điểm đảo chiều}: vì $\hat r \approx 0$, dự báo $\hat P_{t+1} \approx P_t$ luôn đi sau giá thật
        một phiên (Hình~\ref{fig:pred-amzn}b). Khi giá đổi chiều, mô hình chỉ ``thấy'' sau khi việc đã xảy ra.
  \item \textbf{Phân phối đổi theo thời gian}: độ lệch chuẩn của return là \SL{amzn/ret_std/train} trên train nhưng chỉ
        \SL{amzn/ret_std/test} trên test (Chương~\ref{ch:dulieu}), và những phiên lớn nhất đến từ sự kiện ngoài dữ liệu (báo cáo tài chính, tin vĩ mô) mà
        tám đặc trưng kỹ thuật không chứa.
\end{itemize}

Chỉ số $R^2$ trên giá (khoảng 0,99) và biểu đồ ``đường dự báo bám sát đường thật'' là hai cách trình bày dễ làm kết quả
trông tốt hơn thực tế. Với chuỗi giá, cần so với mốc ngây thơ và đo trên return.
